\scanpage{296}
\mspage{449}

Let us return to the general situation of a form $G$ of $\unsure{\mathrm{Gl}(2)_S|\mathrm{GL}(2)_S}$, and suppose that it can be written in the form
\[
G=\underline{\mathrm{Aut}}_{\mathcal O_S\text{-Mod}}(F), \tag{17}
\]
where $F$ is an $\mathcal O_S$-Module \unsure{local\textsuperscript{ly}} free of rank 2,

This is the case as soon as we have a section $x$ of $\unsure{\mathbb P|P}$ over $S$, i.e. a Borel additive subgroup $L$ of $G$, or of $SG$, (in fact we can then take for $F$ the extension of $\mathcal O_S$ by $\mathcal N$ which determines the affine-line bundle $\unsure{\mathbb P|P}\setminus\{x\}$). But \unsure{the|the} $\unsure{\mathbb P|P}$ is realized as \deleted{\illegible{1}}
\[
\unsure{\mathbb P|P}=\unsure{\mathbf P|P}(F)\simeq\unsure{\check{\mathbf P}|\check P}(F) \tag{18}
\]
the projective bundle associated with $F$, \deleted{\unsure{Thus}} \inserted{\unsure{i.e.}} the sections of $\unsure{\mathbb P|P}$ are identified with the ``lines'' $D\subset F$.

We will have
\[
M=\underline{\mathrm{End}}_{\mathcal O_S\text{-mod}}(F)\deleted{\unsure{\simeq}}\simeq\check F\otimes_{\mathcal O_S}F \tag{19}
\]
and the sub-\deleted{bundle}\inserted{Module} \deleted{vector}\inserted{$\mathcal N$} associated with the section $x_D$ of $\unsure{\mathbb P|P}$ \deleted{\unsure{is}} is none other than \deleted{\illegible{1}}
\[
\mathcal N=D^\circ\otimes D
\qquad\bigl(D^\circ\simeq(F/D)^\vee\subset\check F,\text{ the orthogonal to }D\text{ in }\check F\bigr) \tag{20}
\]
its sections are the $u:F\longrightarrow F$ such that
\[
\deleted{u(F)\subset D,\quad u(D)=0,}
\]
\[
u(F)\subset D,\quad u(D)=\{0\}. \tag{21}
\]
\mspage{295}

\scanpage{297}

\uline{Proposition} \deleted{Under the preceding conditions} \inserted{Let $G$ be a form of $\unsure{\mathrm{Gl}(2)_S|\mathrm{GL}(2)_S}$, $\rho,\sigma$ two sections of $G$ such that $\widetilde\rho\widetilde\sigma^{-1}=\widetilde\varepsilon_1$ is a \inserted{\illegible{1}} \unsure{unipotent}-\deleted{\unsure{regular}} section of $PG$, \unsure{linked} \unsure{by the} relation (16). Then} the following conditions are equivalent

(i) $\varepsilon_1$ is \deleted{\illegible{2}} a regular \inserted{section} (i.e. \deleted{\illegible{1}} \inserted{$\ne1$} on each fiber), i.e. $\varepsilon_1=1+N_1$, with $N_1$ a \uline{basis} of $\mathcal N_1$) and $\widetilde\rho^{-1}(x_1)$ is a section of $\unsure{\mathbb P|P}$ \deleted{everywhere} $\ne x_1$ in each fiber.

\margin{\uline{N.B.} We have denoted by $x_1$ the section of $\unsure{\mathbb P|P}$ over $S$ corresponding to $L_1$.}

(ii) \unsure{Ditto|Item}, with $\widetilde\sigma^{-1}$ in place of $\widetilde\rho$.

\margin{$\begin{array}{l}(i\text{ bis})\\(ii\text{ bis})\end{array}\ \big|$ \unsure{Ditto|Item}, with $\widetilde\sigma,\widetilde\rho$ in place of $\widetilde\sigma^{-1},\widetilde\rho^{-1}$.}

(iii) We have (42), i.e. $c=\lambda_0\mathrm{Tr}\,\rho+\mathrm{Tr}\,\sigma$ is \unsure{invertible}.

\uline{Proof.} We have, thanks to the fact that $\widetilde\rho\widetilde\sigma^{-1}\in\Gamma\widetilde L_1$, normalizes $\widetilde L_1$, \inserted{\unsure{or what amounts to the} same, $L_1$,} the relation
\[
\widetilde\rho^{-1}(L_1)=\widetilde\sigma^{-1}(L_1)
\]
or what amounts to the same, we have
\[
\widetilde\rho^{-1}(x_1)=\deleted{\unsure{\widetilde\rho^{-1}}}\inserted{\widetilde\sigma^{-1}}(x_1), \tag{\unsure{22|24}}
\]
which proves the equivalence of (i) and (ii).

Let us show that $(i)\Longleftrightarrow(iii)$. We may assume that $S$ is the spectrum of an algebraically closed field, if it suits us, and we may assume $\lambda_0=1$ (at the cost \deleted{\illegible{1}} of replacing $\rho$ by $\lambda_0\rho$). If we had $\varepsilon_1=1$ i.e. \deleted{$\rho+\sigma$} $1=-\rho\sigma^{-1}$ i.e. $\rho+\sigma=0$, we would have \deleted{$\unsure{\mathrm{Tr}}$} $\mathrm{Tr}(\rho+\sigma)=\mathrm{Tr}\,\rho+\mathrm{Tr}\,\sigma=0$, so (iii) implies the regularity of $\varepsilon_1$ i.e. of $N_1$, \deleted{\illegible{1}} \inserted{we may assume} therefore that $\varepsilon_1$ is regular. We have to show that if we have such a regular unipotent \deleted{\unsure{operator}} \inserted{\unsure{endomorphism}} $\varepsilon_1$\inserted{$=1+N_1$}, of $F$, associated with a $D\subset F$,
\mspage{296}

\scanpage{298}
\mspage{450}

and \unsure{an} \deleted{\unsure{operator} \illegible{1}} \inserted{\unsure{endomorphism}} $\sigma$ of $F$, whence
$\rho=-\varepsilon_1\sigma$, then \deleted{\illegible{2}}
$c=\operatorname{Tr}\sigma+\operatorname{Tr}\rho=\operatorname{Tr}\sigma-\operatorname{Tr}\varepsilon_1\sigma
=-\operatorname{Tr}N_1\sigma$ is invertible iff $D+\sigma^{-1}.D=F$
i.e. $\sigma D+D=F$, \deleted{\unsure{iff}} Let us take a basis
$e_0,e_1$ of $F$ such that $e_1$ generates $D$, and $N_1$ is written as the matrix
$\begin{pmatrix}0&0\\1&0\end{pmatrix}$,
let $\sigma=\begin{pmatrix}a&b\\c&d\end{pmatrix}$ we have
$N_1\sigma=\begin{pmatrix}0&0\\a&b\end{pmatrix}$,
$\operatorname{Tr}N_1\sigma=b$,
$\sigma e_{\unsure{1}}=a e_0+b e_1$, and
$D+\sigma D=D+b\mathcal O_S e_{\unsure{2|0}}$
is equal to $F=\mathcal O_S^2$ iff $b$ invertible, q.e.d.

Let us suppose these equivalent conditions satisfied, \deleted{\unsure{in} $\unsure{P}$} which are also expressed by saying that $L_1$ and
\[
L_0=\widetilde\rho^{-1}(L_1)=\unsure{\widetilde\sigma}^{-1}(L_1)
\tag{23}
\]
are distinct on each fiber, i.e. \unsure{define} Borels $B_1,B_0\subset G$ such that
\[
T=B_1\cap B_0
\tag{24}
\]
is a torus (maximal) of $G$, \deleted{\unsure{which}}
The latter is also the \deleted{\illegible{1}} intersection of the stabilizers $B_0,B_1$ in $G$ of the \inserted{sections}
$x_0,x_1$ of $\unsure{\mathcal P|P}$ which correspond to $L_0,L_1$
(satisfying
\[
x_1=\widetilde\rho x_0=\widetilde\sigma x_0
\tag{25}
\]
), and is \unsure{the inverse} \unsure{image} \unsure{of the subgroup|of the subgroups} (\unsure{privileged|privileged} \inserted{\unsure{toric|toric}}

\mspage{297}

\scanpage{299}

\unsure{called|whose} $\widetilde T$ in $PG=\underline{\operatorname{Aut}}_S(\unsure{\mathcal P|P})$.
We can also say that the datum of two distinct sections \deleted{\unsure{points}}
$x_0,x_1$ of the projective bundle $\unsure{\mathcal P|P}$ \deleted{\unsure{makes it possible to}} amounts to that of the Borel
$B_1$ corresponding to the choice of $x_1$), \uline{and} of the maximal torus $T$ of $B_1$
(which is then written
\[
B_1=T.L_1 \qquad (\unsure{produit}\ \tfrac12\ \text{direct})
\tag{26}
\]

\margin{the pair $(B_1,T)$ is called a \uline{\unsure{Killing} pair}.}

which allows us \unsure{then} to recover $x_0$ as \unsure{one}
\inserted{corresponding to} \unsure{the unique} Borel $B_0$ such that
$T=B_0\cap B_1$.
(The Borels $B_0$ such that $B_0\ne B_1$ on each fiber,
\inserted{i.e. such that $B_0\cap B_1$ is a torus,}
correspond one-to-one to the maximal tori $T$ of $B_1$, by the
\unsure{preceding} formula $T=B_0\cap B_1$.)

This also amounts to: the datum, in addition to $B_1$, \inserted{i.e. to $x_1$,}
of a section \inserted{$x_0$,} of the affine-line bundle
\inserted{$U_1=$}$\unsure{\mathcal P|P}\smallsetminus\{x_1\}$,
\deleted{which \unsure{then makes it possible} to identify $U_1$}
\inserted{i.e. of an \unsure{isomorphism} from $U_1$ onto the}
vector bundle $L_1$ of its translations
\[
U_1\overset{\mathrm{dfn}}{=}\unsure{\mathcal P|P}\smallsetminus\{x_1\}\xrightarrow{\sim}L_1.
\tag{27}
\]

To finish ``pinning'' the situation (i.e.

\mspage{298}

\scanpage{300}
\mspage{451}

reducing $\unsure{\mathcal P|P}$ to the standard projective bundle $\mathbb P^1_S$)
we still have to give ourselves a basis of $L_1$ i.e. a section $x_\infty$ of
\[
L_1^*\simeq U_1\smallsetminus\{x_0\}\simeq\unsure{\mathcal P|P}\smallsetminus\{x_1,x_0\}.
\tag{28}
\]

Now the section $\varepsilon_1$ of $L_1$ is precisely a section of $L_1^*$ -- so we have a complete pinning of the situation.
In terms of $(F,D=D_1)$ corresponding to $L_1$,

\deleted{we have
\[
F=D_1\oplus\underbrace{\sigma^{-1}(D_1)}_{D_0}
=D_1\oplus\underbrace{\rho^{-1}(D_1)}_{D_0},
\]
\unsure{we may assume} (at the cost of tensoring with $D_1^{\otimes-1}$) that $D_1\simeq\mathcal O_S$,
\inserted{\unsure{we have} $N_1|D_0:$}
whence \illegible{1}
\[
\left\{
\begin{aligned}
F&=\mathcal O_S^2\unsure{x}=\mathcal O_S e_1+\mathcal O_S e_0=\mathcal O_S e_0+\mathcal O_S e_1.\\
D_1&=\mathcal O_S e_1,\quad D_0=\mathcal O_S e_0\\
&\unsure{OPS}
\end{aligned}\right.
\tag{\unsure{29}}
\]
we will then have $N_1=(\,\unsure{1}$}
we will have
\[
\left\{
\begin{aligned}
\rho^{-1}(D_1)&=\unsure{\sigma}^{-1}(D_1)\overset{\mathrm{dfn}}{=}D_0\\
F&=D_0\oplus D_1
\end{aligned}\right.
\tag{29}
\]
\deleted{(\unsure{30})}
\[
N_1|D_0:\quad D_0\xrightarrow{\sim}D_1
\tag{30}
\]
and (at the cost of tensoring with $D_1^{\otimes-1}$) \unsure{we may assume}
$D_1\simeq\mathcal O_S$, whence
\[
\left\{
\begin{aligned}
F&\simeq\mathcal O_S^2=\mathcal O_S e_0+\mathcal O_S e_1\\
D_0&=\mathcal O_S e_0,\quad D_1=\mathcal O_S e_1,\\
N_1e_1&=0,\quad N_1e_0=e_1,
\quad\deleted{\text{i.e. }N_1=\begin{pmatrix}0&0\\1&0\end{pmatrix}}
\end{aligned}\right.
\tag{31}
\]

\mspage{299}

\scanpage{301}

i.e.
\[
N_1=\begin{pmatrix}0&0\\1&0\end{pmatrix},\qquad
\varepsilon_1=\begin{pmatrix}1&0\\1&1\end{pmatrix}.
\tag{32}
\]

On the other hand \deleted{\unsure{we will have} $\rho D_0=\sigma\unsure{D}$}
\[
\rho D_0\unsure{=|\ne}\sigma D_0=D_1
\]
\deleted{which is expressed} \inserted{whence}
\[
\rho\sigma^{-1}(D_1)=D_1
\]
\deleted{and we suppose $\rho,\sigma$ chosen over
$\widetilde\rho,\widetilde\sigma$ in such a way that we have
$\rho\sigma^{-1}|D_1=-1$, which (taking into account that
$\widetilde\rho\widetilde\sigma^{-1}=\widetilde\varepsilon_1$)
implies that $\rho\sigma^{-1}(F/D_1)=-1$, i.e.}
\inserted{We \unsure{choose} $\rho,\sigma$ sections in $\mathrm{Gl}(2)$)
over $\widetilde\rho,\widetilde\sigma$
(sections of $\mathrm{GP}(1)\simeq\mathrm{Gl}(2)/\mathrm G_m$
\unsure{satisfying} $\widetilde\rho\widetilde\sigma^{-1}=\widetilde\varepsilon_1$,}
in such a way as to have
\[
\boxed{\rho\sigma^{-1}=-\varepsilon_1}
\tag{33}
\]
which determines each of the two sections $\rho,\sigma$ as a function of the other
(\unsure{=|itself} determined modulo a scalar,
in \deleted{function} \inserted{\unsure{terms}}
of $\unsure{\widetilde\sigma}$ resp. \inserted{of} $\widetilde\rho$).
As
\[
\rho D_0=D_1,
\]
we will have
\[
\rho=\begin{pmatrix}0&\unsure{u|U}\\\unsure{v|V}&\unsure{w|W}\end{pmatrix},
\qquad \det\rho=-\unsure{uv|UV}\ \text{invertible}
\tag{*}
\]
and we can remove the scalar \unsure{ambiguity|ambiguity} in $\rho$
by \deleted{\unsure{writing}} normalizing by
\[
\rho=\begin{pmatrix}0&1\\\unsure{v|V}&\unsure{w|W}\end{pmatrix}
\qquad \unsure{v|V}\in\Gamma\mathrm G_m.
\tag{34}
\]

\mspage{300}